\section{Introduction}
\label{sec:intro}

An agent that interacts with a user over weeks accumulates a long history of conversations, and a growing part of this history consists of shared images such as screenshots, photos, product pictures, and charts. When the user later asks a question, the memory of the agent must return the past evidence that answers it. Recent controlled studies show that a simple baseline is hard to beat for this task: keeping the raw history and retrieving it by embedding similarity matches or exceeds memory systems that extract, rewrite, or reorganize the history before any question arrives~\citep{bao2026rawrag,zhou2026explicit,sharma2026selection,peng2026memorytax}.

Similarity retrieval rests on an implicit assumption, namely that the evidence for a question consists of the items that are individually most similar to it. The retriever therefore scores every memory item independently and returns the $k$ nearest neighbors of the question. On the text benchmarks where the observation was reported, this assumption is adequate once the budget is large: in our replication, cosine retrieval of raw turns is statistically indistinguishable from reading the entire history on LoCoMo~\citep{maharana2024locomo} and LongMemEval~\citep{wu2025longmemeval} (Section~\ref{sec:prelim}).

Multimodal memory strains this assumption in a specific way. Many questions refer to a \emph{set} of past items rather than to a single item, such as the five advertisements of one brand, the summary cards of day 9 and day 11, or every checkpoint of a night route. The answer depends on every member, so a single missing image can change a count or a comparison. The members of such a set are often named by narrative aliases, such as ``the D9 summary'' or ``the month-wrap comparison trend'', which appear neither in the dialogue around the image nor in its caption and can be resolved only from the place of the item in the user history. Top-$k$ ranking has no notion of set completeness, and its probability of returning the whole set decays sharply with the size of the set (Figure~\ref{fig:setsize}).

A memory that retrieves sets faces four challenges. First, retrieval must be measured per set, because recall@$k$ and coverage average over items and hide incomplete sets. Second, the retriever must decide membership for items whose relation to the question is defined by the structure of the history, with a budget that follows the size of the set. Third, better sets must improve the answers on data that played no role in the design, where an index may hide the evidence and a lightweight reader must aggregate several members. Fourth, a retriever that reasons over the whole history at question time must remain affordable, since it adds a model call to every question.

To address the first two challenges, we score retrieval by full-set recall, the per-question form of the all-recall measure used in multi-hop retrieval~\citep{gutierrez2024hipporag}, and we study query-time set selection (\qss{}), which separates navigation from perception (Figure~\ref{fig:overview}). Memory content stays raw. At ingestion, the memory writes one short index line per round with the date, a dialogue snippet, and a one-sentence description of each image; at question time, a language model reads the whole index and returns a set of rounds whose size it decides, and the reader answers from the original images of these rounds. Reading the whole index allows the selector to resolve an alias against the place of an item in the history, which a similarity score between the question and the item cannot do, and returning a set allows the budget to follow the number of members. \qss{} belongs to the family of index-based navigation methods, such as gist memory with look-up~\citep{lee2024readagent} and summary-tree navigation~\citep{chen2023memwalker}, and differs from them in returning an unordered set with model-chosen cardinality. We isolate what this set formulation adds with controls that use the same model and candidates, namely listwise ranking cut at the same size, with and without the recall-oriented instruction of \qss{}, and graded relevance with a dynamic cut-off. For the third challenge, we test the method, which we developed on MemEye~\citep{guo2026memeye}, on MemLens~\citep{ren2026memlens}, a benchmark from other authors whose visual necessity was validated by image ablation, with hypotheses and conditions registered before any run.

\begin{figure}[t]
\centering
\begin{tikzpicture}[font=\footnotesize, node distance=4mm and 9mm,
  box/.style={draw, rounded corners=2pt, align=center, inner sep=3pt, fill=white, minimum height=1.25cm},
  arr/.style={-{Latex[length=2mm]}, thick}]
\node[box, minimum width=2.7cm] (hist) {\textbf{Raw memory}\\rounds: dialogue +\\original images};
\node[box, right=of hist, minimum width=3.0cm] (idx) {\textbf{Memory index}\\one line per round:\\date, snippet, image gist};
\node[box, right=of idx, minimum width=3.3cm, fill=red!6] (sel) {\textbf{Query-time selection}\\an LM reads the index\\and the question};
\node[box, right=of sel, minimum width=2.5cm] (read) {\textbf{Reader}\\answers from the\\raw rounds in $S$};
\draw[arr] (hist) -- node[above, font=\scriptsize]{write} (idx);
\draw[arr] (idx) -- (sel);
\draw[arr] (sel) -- node[above, font=\scriptsize]{$S$} (read);
\draw[arr, dashed] (hist.south) -- ++(0,-0.4) -| node[pos=0.25, below, font=\scriptsize]{the reader opens the selected rounds with their original images} (read.south);
\draw[arr, dotted] (sel.north) -- ++(0,0.35) -| node[pos=0.25, above, font=\scriptsize]{resolution write-back: question and selected set} (idx.north);
\end{tikzpicture}
\caption{Query-time set selection (\qss{}). The index is written once per round and serves only navigation; the selector decides set membership against the whole index at question time, and the reader sees the original images of the selected rounds. Resolution write-back appends answered questions and their selected sets to the index.}
\label{fig:overview}
\end{figure}

On MemEye, a strong selector returns complete evidence sets for 90.3\% of questions while opening 7.3 rounds, keeps this recall across set sizes, and enlarges its budget when the set is larger. The controls show that most of this recall comes from model judgment and from the recall-oriented instruction, and that the unordered set adds a small further gain and decides how many rounds to open. The held-out test bounds these findings. On MemLens, \qss{} again doubles the full-set recall of equal-size cosine retrieval, but its answers only match those of top-40 retrieval with a lightweight reader and fall 0.080 below them with a strong reader. A simple measurement accounts for much of the contrast: almost all question words that occur in MemEye evidence rounds are visible in their index lines, whereas most of them are hidden in MemLens, where the evidence is an aside inside long messages about other topics. Neither write-time organization by the same strong model, nor prompting interventions, nor selectors from other model families with the same prompt reproduce the selection quality of the developed selector.

We then turned this diagnosis into mechanisms for the last two challenges. A failure analysis shows that lightweight readers fail to aggregate evidence sets even when they are complete: given the full annotated evidence, the lightweight reader answers 6\% of the MemLens arithmetic questions, against 67\% for the strong reader. Visibility-gated set reading (VGS) therefore estimates index visibility at question time and, for low-visibility questions, widens the index lines and passes the reader per-member notes written by a lightweight model, which make the members of the set and their values explicit. Developed on MemEye and a third of MemLens, VGS improves the lightweight reader on the remaining MemLens questions by $+0.071$ $[+0.032, +0.109]$ over top-40 retrieval in a registered test, whereas making sets more complete without such help lowers accuracy; on the external WorldMemArena benchmark, however, the frozen design gains no significant accuracy although its visibility is as low as that of MemLens. To keep the selection call affordable, we exploit the fact that the index of a persistent memory changes only by appending: a compact index is stored once per user in an explicit context cache, and a router that reads only the question sends single-item questions to retrieval. With these mechanisms, a strong reader with set selection is more accurate than top-40 retrieval on MemEye at 64\% of its cost, and the compact index halves the cost of VGS on held-out questions without lowering its accuracy; the cost advantage disappears when each cached index serves only a few questions.

The contributions of this paper are as follows. First, we confirm that raw similarity retrieval suffices on text benchmarks when the budget is large, and we show that its limit in multimodal memory is set completeness, which full-set recall measures. Second, we study query-time set selection over a lightweight index and decompose its recall, with same-model controls on two benchmarks, into model judgment, a recall-oriented instruction, and a set output of model-chosen size; the capability did not transfer to the lightweight selectors, write-time organizations, and cross-family selectors that we tested. Third, we bound these claims with a pre-registered held-out evaluation and identify index visibility as a measurable predictor of when set selection improves answers. Fourth, we propose visibility-gated set reading, which acts on this predictor and on the aggregation failure of lightweight readers; it improves them in registered tests on held-out MemLens data but not on an external benchmark, which bounds the role of visibility as a predictor. Fifth, we show with measured costs that a compact, cache-native index and set-aware routing bring set selection close to the cost of retrieval, and below it for a strong reader on a persistent memory whose cached index serves enough questions. Throughout, we report the negative results that constrain designs, including aggressive set completion, write-time groups and aliases, cheaper selection calls, and reader-driven zooming.
